%% BLOCK 1 -- Introduction to Machine Learning.
%% Fragment of ai_foundations.tex; not a standalone document.
%% Source: developers.google.com/machine-learning/intro-to-ml, all three pages
%% (What is Machine Learning?, Supervised Learning, Test Your Understanding).

\actpage{1}{Introduction to Machine Learning}{What the four kinds of ML system are, and what the five parts of a supervised one do.}

% ======================================================================
\sectionsubtitle{Four categories, by how the system learns.}
\section[Types of ML]{Types of ML Systems}

\begin{frame}{Types of ML systems}
  \begin{withmargin}
    \begin{tdmain}
      ML systems fall into one or more of the following categories based on how
      they learn to make predictions or generate content.
      \vskip0.6em
      \begin{itemize}
        \item Supervised learning
        \item Unsupervised learning
        \item Reinforcement learning
        \item Generative AI
      \end{itemize}
      \vskip0.6em
      \footnotesize
      The categories are not exclusive. Generative models are initially trained
      using an \gterm{unsupervised} approach, then sometimes trained further
      using supervised or reinforcement learning.
    \end{tdmain}
    \begin{tdmargin}
      {\color{tdfg}\textbf{Course length: 20 minutes.}}
      \par\medskip
      This course introduces ML concepts. It does not cover how to implement ML
      or work with data.
      \par\medskip
      Objectives: the different types of machine learning, the key concepts of
      supervised ML, and how solving problems with ML differs from traditional
      approaches.
    \end{tdmargin}
  \end{withmargin}
  \slidesources{GoogleIntroML2025}
\end{frame}

\begin{frame}{Supervised learning}
  \begin{withmargin}
    \begin{tdmain}
      Supervised learning models can make predictions after seeing lots of data
      with the correct answers and then discovering the connections between the
      elements in the data that produce the correct answers.
      \vskip0.9em
      These ML systems are \alertbf{supervised} in the sense that a human gives
      the ML system data with the known correct results. Two of the most common
      use cases are regression and classification.
    \end{tdmain}
    \begin{tdmargin}
      {\color{tdfg}\textbf{The old-exams analogy.}}
      \par\medskip
      This is like a student learning new material by studying old exams that
      contain both questions and answers. Once the student has trained on enough
      old exams, the student is well prepared to take a new exam.
    \end{tdmargin}
  \end{withmargin}
\end{frame}

\begin{frame}{Regression predicts a number}
  \begin{withmargin}
    \begin{tdmain}
      A regression model predicts a \gterm{numeric value}. A weather model that
      predicts the amount of rain, in inches or millimeters, is a regression
      model.
      \vskip0.7em
      \scriptsize
      \begin{tabular}{@{}p{0.20\linewidth}p{0.44\linewidth}p{0.28\linewidth}@{}}
        \toprule
        \textbf{Scenario} & \textbf{Possible input data} & \textbf{Numeric prediction}\\
        \midrule
        Future house price & Square footage, zip code, number of bedrooms and
          bathrooms, lot size, mortgage interest rate, property tax rate,
          construction costs, number of homes for sale in the area
          & The price of the home\\[0.4em]
        Future ride time & Historical traffic conditions, distance from
          destination, weather conditions
          & The time to arrive at a destination\\
        \bottomrule
      \end{tabular}
    \end{tdmain}
    \begin{tdmargin}
      Historical traffic conditions are gathered from smartphones, traffic
      sensors, ride-hailing and other navigation applications.
      \par\medskip
      Energy usage is measured in kilowatt-hours, which is a number, so a model
      predicting energy usage for commercial buildings is a regression model.
    \end{tdmargin}
  \end{withmargin}
\end{frame}

\begin{frame}{Classification predicts a category}
  \begin{withmargin}
    \begin{tdmain}
      Classification models predict the likelihood that something belongs to a
      category. Unlike regression models, whose output is a number,
      classification models output a value that states whether or not something
      belongs to a particular category. Classification models are used to
      predict if an email is spam or if a photo contains a cat.
      \vskip0.9em
      \textbf{Binary} classification models output a value from a class that
      contains only two values, for example \texttt{rain} or \texttt{no rain}.
      \textbf{Multiclass} classification models output a value from a class that
      contains more than two values, for example \gterm{\texttt{rain}},
      \gterm{\texttt{hail}}, \gterm{\texttt{snow}}, or \gterm{\texttt{sleet}}.
    \end{tdmain}
    \begin{tdmargin}
      {\color{tdfg}\textbf{Turning a number into classes.}}
      \par\medskip
      To classify an airplane ticket cost as high, average, or low: find the
      average cost from departure to destination airport, determine the
      thresholds that constitute high, average, and low, then compare the
      predicted cost to the thresholds and output the category the value falls
      within.
    \end{tdmargin}
  \end{withmargin}
\end{frame}

\begin{frame}{Unsupervised learning}
  \begin{withmargin}
    \begin{tdmain}
      An unsupervised learning model aims to identify meaningful patterns in a
      dataset. Many unsupervised learning models rely on a technique called
      \gterm{clustering} to organize similar data into groups.
      \vskip0.9em
      Clustering differs from classification because the categories aren't
      defined by you. An unsupervised model might cluster a weather dataset
      based on temperature, revealing segmentations that define the seasons. You
      might then attempt to name those clusters based on your understanding of
      the dataset.
    \end{tdmain}
    \begin{tdmargin}
      {\color{tdfg}\textbf{The distinguishing question.}}
      \par\medskip
      A supervised approach is given data that contains the correct answer. Its
      job is to find connections in the data that produce the correct answer.
      \par\smallskip
      An unsupervised approach is given data without the correct answer. Its job
      is to find groupings in the data. It does not know what the clusters mean.
    \end{tdmargin}
  \end{withmargin}
\end{frame}

\begin{frame}{Reinforcement learning}
  \begin{withmargin}
    \begin{tdmain}
      Reinforcement learning models make predictions by getting \gterm{rewards}
      or penalties based on actions performed within an environment. A
      reinforcement learning system generates a policy that defines the best
      strategy for getting the most rewards.
    \end{tdmain}
    \begin{tdmargin}
      Reinforcement learning is used to train robots to perform tasks, like
      walking around a room, and software programs like AlphaGo to play the game
      of Go.
    \end{tdmargin}
  \end{withmargin}
\end{frame}

\begin{frame}{Generative AI}
  \begin{withmargin}
    \begin{tdmain}
      Generative AI is a class of models that creates content from user input.
      It can create unique images, music compositions, and jokes; it can
      summarize articles, explain how to perform a task, or edit a photo.
      \vskip0.7em
      We discuss generative models by their inputs and outputs, typically
      written as \gterm{input type} to \gterm{output type}:
      \vskip0.4em
      \footnotesize
      \begin{itemize}
        \item Text-to-text, text-to-image, text-to-video
        \item Text-to-code, text-to-speech
        \item Image and text-to-image
      \end{itemize}
    \end{tdmain}
    \begin{tdmargin}
      A model can take an image as input and create an image and text as output,
      or take an image and text as input and create a video as output.
      \par\medskip
      Generative models help businesses refine ecommerce product images, by
      removing distracting backgrounds or raising the quality of low-resolution
      images.
    \end{tdmargin}
  \end{withmargin}
\end{frame}

\begin{frame}{How generative AI works}
  \begin{withmargin}
    \begin{tdmain}
      At a high level, generative models learn patterns in data with the goal to
      produce new but similar data. Generative models are like:
      \vskip0.5em
      \begin{itemize}
        \item Comedians who learn to imitate others by observing people's
          behaviors and style of speaking
        \item Artists who learn to paint in a particular style by studying lots
          of paintings in that style
        \item Cover bands that learn to sound like a specific music group by
          listening to lots of music by that group
      \end{itemize}
    \end{tdmain}
    \begin{tdmargin}
      To produce unique and creative outputs, generative models are initially
      trained using an unsupervised approach, where the model learns to
      \gterm{mimic} the data it's trained on.
      \par\medskip
      The model is sometimes trained further using supervised or reinforcement
      learning on data for tasks it might be asked to perform, for example
      summarize an article or edit a photo.
    \end{tdmargin}
  \end{withmargin}
  \keyterms{classification model; clustering; model; policy; prediction;
    regression model; reinforcement learning; reward; supervised learning;
    training; unsupervised learning}
\end{frame}

% ======================================================================
\sectionsubtitle{Data, model, training, evaluating, inference.}
\section[Supervised concepts]{Foundational Supervised Learning Concepts}

\begin{frame}{Data is the driving force of ML}
  \begin{withmargin}
    \begin{tdmain}
      Data comes in the form of words and numbers stored in tables, or as the
      values of pixels and waveforms captured in images and audio files. We
      store related data in \gterm{datasets}.
      \vskip0.9em
      Datasets are made up of individual \textbf{examples} that contain
      \textbf{features} and a \textbf{label}. You could think of an example as
      analogous to a single row in a spreadsheet. Features are the values that a
      supervised model uses to predict the label. The label is the answer, or
      the value we want the model to predict.
    \end{tdmain}
    \begin{tdmargin}
      {\color{tdfg}\textbf{A weather model.}}
      \par\medskip
      Features: latitude, longitude, temperature, humidity, cloud coverage, wind
      direction, atmospheric pressure. Label: rainfall amount.
      \par\medskip
      Examples that contain both features and a label are called
      \textbf{labeled examples}. \textbf{Unlabeled examples} contain features
      but no label; after you create a model, the model predicts the label from
      the features.
    \end{tdmargin}
  \end{withmargin}
\end{frame}

\begin{frame}{Size and diversity}
  \begin{withmargin}
    \begin{tdmain}
      A dataset is characterized by its size and diversity. \textbf{Size}
      indicates the number of examples. \textbf{Diversity} indicates the range
      those examples cover. Good datasets are \alertbf{both large and highly
      diverse}.
      \vskip0.9em
      A large dataset doesn't guarantee sufficient diversity, and a dataset that
      is highly diverse doesn't guarantee sufficient examples.
    \end{tdmain}
    \begin{tdmargin}
      A dataset might contain 100 years worth of data, but only for the month of
      July. Using it to predict rainfall in January would produce poor
      predictions.
      \par\smallskip
      Conversely, a dataset might cover only a few years but contain every
      month. It might produce poor predictions because it doesn't contain enough
      years to account for variability.
    \end{tdmargin}
  \end{withmargin}
\end{frame}

\begin{frame}{More features is not automatically better}
  \begin{withmargin}
    \begin{tdmain}
      A dataset can also be characterized by the number of its features. Some
      weather datasets might contain hundreds of features, ranging from
      satellite imagery to cloud coverage values. Others might contain only
      three or four, like humidity, atmospheric pressure, and temperature.
      \vskip0.9em
      Datasets with more features can help a model discover additional patterns
      and make better predictions. However, datasets with more features don't
      always produce models that make better predictions, because some features
      might have \alertbf{no causal relationship} to the label.
    \end{tdmain}
    \begin{tdmargin}
      {\color{tdfg}\textbf{Predictive power.}}
      \par\medskip
      In a used-car dataset with price as the label, make/model, year, and miles
      are likely among the strongest predictors. Color, height, tire size, and
      wheel base are not.
    \end{tdmargin}
  \end{withmargin}
\end{frame}

\begin{frame}{Model}
  \vfill
  \begin{takeaway}
    A model is the complex collection of numbers\\
    that define the mathematical relationship\\
    from input feature patterns to output label values.
  \end{takeaway}
  \vfill
  {\small The model discovers these patterns through training.}
\end{frame}

\begin{frame}{Training}
  \begin{withmargin}
    \begin{tdmain}
      To train a model, we give the model a dataset with labeled examples. The
      model's goal is to work out the best solution for predicting the labels
      from the features. The model finds the best solution by comparing its
      predicted value to the label's actual value. Based on the difference
      between the predicted and actual values, defined as the \gterm{loss}, the
      model gradually updates its solution.
      \vskip0.6em
      \footnotesize
      \begin{itemize}
        \item The model takes in a single labeled example and provides a prediction.
        \item It compares its predicted value with the actual value and updates
          its solution.
        \item It repeats this for each labeled example in the dataset.
      \end{itemize}
    \end{tdmain}
    \begin{tdmargin}
      If the model predicted \texttt{1.15 inches} of rain, but the actual value
      was \texttt{.75 inches}, the model modifies its solution so its prediction
      is closer to \texttt{.75 inches}.
      \par\medskip
      After the model has looked at each example in the dataset, in some cases
      multiple times, it arrives at a solution that makes the best predictions,
      on average, for each of the examples.
    \end{tdmargin}
  \end{withmargin}
\end{frame}

\begin{frame}{Why large and diverse datasets produce a better model}
  \begin{withmargin}
    \begin{tdmain}
      The model gradually learns the correct relationship between the features
      and the label. This \alertbf{gradual} understanding is also why large and
      diverse datasets produce a better model: the model has seen more data with
      a wider range of values and has refined its understanding of the
      relationship between the features and the label.
      \vskip0.9em
      During training, ML practitioners can make subtle adjustments to the
      configurations and features the model uses. Certain features have more
      predictive power than others, so practitioners can select which features
      the model uses during training.
    \end{tdmain}
    \begin{tdmargin}
      Suppose a weather dataset contains \texttt{time\_of\_day} as a feature. An
      ML practitioner can add or remove \texttt{time\_of\_day} during training
      to see whether the model makes better predictions with or without it.
    \end{tdmargin}
  \end{withmargin}
\end{frame}

\begin{frame}{Evaluating, then inference}
  \begin{withmargin}
    \begin{tdmain}
      We evaluate a trained model to determine how well it learned. When we
      evaluate a model, we use a labeled dataset, but we only give the model the
      dataset's \textbf{features}. We then compare the model's predictions to
      the label's true values. Depending on the results, we might do more
      training and evaluating before deploying.
      \vskip0.7em
      Once we're satisfied, we can use the model to make predictions, called
      \gterm{inferences}, on unlabeled examples. In the weather app, we give the
      model the current conditions, like temperature, atmospheric pressure, and
      relative humidity, and it predicts the amount of rainfall.
    \end{tdmain}
    \begin{tdmargin}
      {\color{tdfg}\textbf{If the predictions are far off.}}
      \par\medskip
      Retrain using only the features you believe have the strongest predictive
      power for the label.
      \par\smallskip
      Retrain using a larger and more diverse dataset.
      \par\medskip
      Most models require multiple rounds of training.
    \end{tdmargin}
  \end{withmargin}
  \keyterms{example; feature; inference; labeled example; label; loss;
    prediction; training}
\end{frame}
